% ============================================================================
% SECTION 3 — APPROACH: THE TDBF PIPELINE
% Paper budget: ~3 pages + one figure + one tool table. This is the
% contribution section, so it documents the system AS BUILT
% (src/tdbf/agent/*, scripts/07_deployable_pipeline.py) — not a design
% proposal. A paper describes mechanism, then says why the mechanism matters;
% implementation detail that does not change a result goes to the artifact.
% Source to adapt: thesisproposal/main.tex §4.
% ============================================================================
\section{Approach}
\label{sec:approach}

\subsection{Pipeline Overview}
\label{sec:method-overview}
% TODO (~2 paragraphs + Fig. 1). The strict order as implemented:
% BRT agent → verify the test FAILS on the buggy revision (gold-free gate) →
% fix agent → offline scoring. State the invariant the architecture buys:
% the repair stage is unreachable until a failure has been executed and seen.
% Figure: \includegraphics[width=\linewidth]{tdbf_pipeline.pdf}
% (single-column acmart — \linewidth is the full text width; there is no
% figure*/table* distinction to worry about in this layout).

\subsection{The Gold-Free Invariant}
\label{sec:method-goldfree}
% TODO (~2 paragraphs). The design contribution that emerged during the work.
% No agent ever holds a tool that observes the gold patch, the gold-fixed tree,
% or the held-out FAIL_TO_PASS tests; reference data is an offline measurement
% instrument only. Enforced programmatically: allowlisted tool registries plus
% assert_gold_free / assert_fix_gold_free guards, and context_safe_dict() as
% the only instance view an agent receives. Why it matters, in one sentence a
% reviewer can quote: a pipeline whose accept decision depends on the reference
% patch cannot be deployed — this is the difference between an oracle study and
% a system.

\subsection{Agent Runtime and Tools}
\label{sec:method-runtime}
% TODO (~1 paragraph + Table I). Vendor-agnostic tool-use loop (runtime.py):
% terminal submit_* tools, iteration and token caps, persistent per-instance
% container (session.py) so many tool calls pay startup once.
% Table I — tools per agent (this table also defines Arm A by subtraction):
%   BRT agent:  read_file, grep, glob, run_python, run_test (buggy only),
%               submit_test
%   Fix agent:  read_file, grep, glob, edit_file, run_brt, run_tests, revert,
%               submit_fix
% One sentence on the load-bearing detail: run_brt has the validated BRT baked
% in, so the fix agent cannot alter its own success criterion.

\subsection{Phase 1: Test Generation and the Acceptance Gate}
\label{sec:method-phase1}
% TODO (~2 paragraphs). As built: the agent probes the buggy tree, writes the
% test, runs it. The gate accepts iff the test reproduces a CLEAN failure on the
% buggy revision — it reached its assertion and failed there, rather than
% crashing before asserting. Say why the distinction is load-bearing: a crash
% fails identically on fixed code, so it can never be a valid BRT.
% Candidate set = L3 ∪ L4 (Section~\ref{sec:method-levels}); the L3/L4 split is
% precisely what cannot be told without the reference patch, which is why a
% deployable pipeline must attempt to fix L3s too. Mention the bounded
% gate-retry (fresh conversation, gold-free feedback) in one sentence.

\subsection{Phase 2: Fix Generation}
\label{sec:method-phase2}
% TODO (~1–2 paragraphs). As built: the BRT is given to the fix agent as a
% GUIDE, not as the oracle (prompt design); bail-out rule; targeted regression
% testing; revert discipline; the submitted patch is the git diff of the tracked
% tree with the BRT file excluded automatically.

\subsection{The Reactive Baseline}
\label{sec:method-baseline}
% TODO (~1 paragraph). Arm A: same model, tools, budgets, container, and
% scoring; run_brt removed and run_tests kept; no Phase 1; the submit_fix
% description carries no reference to a reproduction test. State the fairness
% rule explicitly — the baseline receives issue text, repository context, and
% the repository's own tests, exactly as the test-first arm does.

\subsection{Outcome Taxonomy and Metrics}
\label{sec:method-levels}
% TODO (~1 paragraph + Table II). The L0–L4 taxonomy of a generated test by its
% buggy-vs-fixed behaviour, mapped to SWE-bench vocabulary (L4 = Fail→Pass =
% plausible BRT; L3 = Fail→Fail = candidate-not-plausible). Define, as display
% math or inline: C2P = L4/(L3+L4); Valid-BRT = L4/N; resolved = all held-out
% FAIL_TO_PASS pass AND PASS_TO_PASS preserved. Say once, here, that L0–L4
% labels are computed offline and never gate anything at run time.
